\documentclass[11pt,letterpaper]{article}
\usepackage[margin=1in]{geometry}
\usepackage[T1]{fontenc}
\usepackage{newtxtext,newtxmath}
\usepackage{microtype}
\usepackage[authoryear,round]{natbib}
\usepackage[hidelinks]{hyperref}

\hypersetup{
  pdftitle={Lateral OFC represents reward identity during conditioning},
  pdfsubject={Proposed Results text for PI review}
}
\setlength{\parindent}{0pt}
\setlength{\parskip}{0.8\baselineskip}
\pagestyle{empty}

\begin{document}

{\small\sffamily Draft for PI review \hfill November 12, 2026\par}
\vspace{0.4\baselineskip}
{\Large\bfseries Lateral OFC represents reward identity during conditioning\par}
\vspace{0.4\baselineskip}

The preferential representation of preconditioned outcomes in lateral OFC may reflect the prioritization of inferred outcome identities when both cue classes share limited representational resources. We therefore asked whether directly conditioned cues would evoke odor-related representations during conditioning, when only conditioned cues were presented and participants explicitly reported which odor would follow each cue. Using the same independently defined lateral and medial OFC regions, we tested whether cues predicting the same food odor evoked more similar neural patterns than cues predicting different odors. Analyses used single-trial activity estimates spanning the combined cue and prediction period (0--5\,s), pooled correct trials across all available conditioning runs ($n = 29$), and excluded clean-air trials. Pattern comparisons were restricted to presentations from different runs.

Neural patterns in lateral OFC showed a small positive correspondence with the odor-identity model (mean Spearman $\rho = 0.0131$, $t(28) = 2.01$, $p = 0.027$), whereas the corresponding effect was not significant in medial OFC (mean Spearman $\rho = 0.0068$, $t(28) = 1.20$, $p = 0.120$; both $p$ values one-sided and uncorrected). These findings provide suggestive evidence of odor-related information in lateral OFC during explicit prediction of directly associated outcomes, consistent with previous findings of identity-specific reward representations in this region \citep{howard2015identity,howard2017identity}.

{\small
\bibliographystyle{plainnat}
\bibliography{../main/sn-bibliography}
}

\end{document}
